\documentclass[10pt,a4paper]{amsart}
\usepackage[margin=19mm]{geometry}
\usepackage{amsmath,amssymb,mathtools}
\usepackage[scr=rsfs,cal=euler]{mathalfa}
\usepackage{xfrac}
\usepackage[unicode,hidelinks,bookmarks=false]{hyperref}
\newcommand{\ie}{i.e.\ }
\newcommand{\cf}{cf.\ }
\newcommand{\resp}{respectively\ }
\newcommand{\Subsection}{\subsection}
\providecommand{\nobreakdash}{\nobreak\hbox{-}}
\setlength{\emergencystretch}{2em}
\multlinegap0pt
\usepackage{amssymb}
\usepackage{bm}
\usepackage{mathtools}
\usepackage{algorithm}\usepackage{algpseudocode}
\usepackage{float}
\newtheorem{proposition}{Proposition}
\newtheorem{theorem}{Theorem}
\title{Higher Degree $t$-Hermitian Forms and Positivity-Preserving Contractions}
\author{Isaac Dobes}
\date{Source: arXiv:2602.21048v3 (CC BY 4.0)}
\begin{document}
\maketitle

\noindent\emph{Source and licence.} Higher Degree $t$-Hermitian Forms and Positivity-Preserving Contractions. Version arXiv:2602.21048v3 (CC BY 4.0), distributed by arXiv under the Creative Commons Attribution 4.0 International licence, http://creativecommons.org/licenses/by/4.0/, as recorded in the arXiv licence metadata for this exact version and shown on the abstract page https://arxiv.org/abs/2602.21048v3. Full licence notice: \texttt{licenses/arxiv-2602.21048-CC-BY-4.0.txt}.\par\smallskip

\noindent\textit{Editorial scope.} Counted unit: the article's theorem theorem (source0.tex lines 856--867). The article's complete original proof of it is delivered with it (source0.tex lines 868--898, one \texttt{proof} environment of 31 lines). No mathematics is written, repaired or re-derived: the statement and the proof are the article's own lines, in the article's own notation, and no retained line is substituted or re-flowed.  Retained uncounted as the source-local context the proof needs: lines 818--820; lines 822--831.  Nothing was omitted from any retained span, so the package carries no omission record.  No retained line contains a citation, so the wrapper carries no bibliography and no bibliography entry.  No retained line contains a figure environment, an image-inclusion command or drawing code, so no image is delivered and the package declares no asset dependency.  Editorial formatting only, in addition: an AMS \texttt{amsart} preamble that restates the article's own declarations and macros used by the retained lines, namely the environments \texttt{proposition}, \texttt{theorem} on separate counters, together with the standard packages those lines need.  The retained environments print in this document as Proposition 1, Theorem 1 in order of appearance, whereas the article prints the same items with the numbering of its own full document.  The complete native source of the article as distributed in the arXiv e-print for this exact version is bundled unchanged as \texttt{sources/195241/source0.tex}.
    \begin{equation}\label{contraction hypermatrix in frequency domain}
        \Phi_c(\mathcal{A}) = \sum\limits_{\ell=1}^p\overline{c[\ell]}\mathcal{A}_{(\ell)} = \sum\limits_{\ell,m=1}^p\overline{c[\ell]}\omega_p^{-(\ell-1)(m-1)}\widehat{\mathcal{A}}_{(m)} = \frac{1}{p}\sum\limits_{m=1}^p\overline{\widehat{c}[m]}\widehat{\mathcal{A}}_{(m)}.
    \end{equation}

\begin{proposition}[Quantitative Lower Bound for Contraction Induced Forms]\label{Quantitative Lower Bound for Contraction Induced Forms}
    Let $\mathcal{A}\in \mathbb{C}^{n\times\dots\times n\times p}$ be an order $2k+1$ $t$-Hermitian positive semidefinite hypermatrix and $c\in \mathbb{T}_p^+$, and define 
    \[\alpha_{\ell} := \min\limits_{\|z\|=1}h_{\widehat{\mathcal{A}}_{(\ell)}}(z).\]
    Then \[h_{\Phi_c(\mathcal{A})}(z) \geq \gamma_c(\mathcal{A})\|z\|^{2k}\qquad \forall z\in \mathbb{C}^n,\]
    where 
    \[\gamma_c(\mathcal{A}) := \frac{1}{p}\sum\limits_{l=1}^p\widehat{c}[\ell]\alpha_{\ell},\]
    hence $h_{\Phi_c(\mathcal{A})}(x)$ is Hermitian positive semidefinite. 
    
    Furthermore, if additionally $\mathcal{A}$ is $t$-Hermitian positive definite and $c\neq 0$, then $\gamma_c(\mathcal{A}) > 0$, hence $h_{\Phi_c(\mathcal{A})}(x)$ is Hermitian positive definite. 
\end{proposition}

\begin{theorem}[Characterization of Positivity-Preserving Contractions]
    Let $c\in \mathbb{T}_p$. The following are equivalent:
    \begin{itemize}
        \item[(i)] $c\in \mathbb{T}_p^+$; that is, $\widehat{c}[\ell]\geq 0$ for each $\ell\in [p]$. 
        \item[(ii)] $\Phi_c$ \textbf{weakly preserves positivity}, meaning that for every order $2k+1$ $t$-Hermitian positive semidefinite hypermatrix $\mathcal{A}\in \mathbb{C}^{n\times\dots\times n\times p}$, $h_{\Phi_c(\mathcal{A})}(x)$ is Hermitian positive semidefinite. 
    \end{itemize}
    Additionally, the following are equivalent:
    \begin{itemize}
        \item[(a)] $c\in \mathbb{T}_p^+\setminus \{0\}$; that is, $\widehat{c}[\ell]\geq 0$ for all $\ell\in [p]$ and $\widehat{c}[\ell_*] > 0$ for atleast one $\ell_*\in [p]$. 
        \item[(b)] $\Phi_c$ \textbf{strongly preserves positivity}, meaning that for every order $2k+1$ $t$-Hermitian positive definite hypermatrix $\mathcal{A}\in \mathbb{C}^{n\times\dots\times n\times p}$, $h_{\Phi_c(\mathcal{A})}(x)$ is Hermitian positive definite. 
    \end{itemize}
\end{theorem}

\begin{proof}
    Both equation \eqref{contraction hypermatrix in frequency domain} and Proposition \ref{Quantitative Lower Bound for Contraction Induced Forms} immediately give $(i)\Rightarrow (ii)$. 

    Conversely, suppose that $\Phi_c$ weakly preserves positivity; that is, for every order $2k+1$ $t$-Hermitian positive semidefinite hypermatrix $\mathcal{A}\in \mathbb{C}^{n\times\dots\times n\times p}$, the $c$-contraction hypermatrix
    \[\Phi_c(\mathcal{A}) = \sum\limits_{\ell=1}^p\overline{c[\ell]}\mathcal{A}_{(\ell)}\]
    represents a classical Hermitian positive semidefinite form. Fix $m\in [p]$, let $\mathcal{B}\in\mathbb{C}^{n\times\dots\times n}$ be any nonzero order $2k$ Hermitian positive semidefinite hypermatrix, and define the order $2k+1$ hypermatrix $\mathcal{A}\in \mathbb{C}^{n\times\dots\times n\times p}$ so that 
    \[\widehat{\mathcal{A}}_{(m)} = \mathcal{B}\]
    and 
    \[\widehat{\mathcal{A}}_{(\ell)} = 0\]
    for all $\ell\neq m$. By construction, every frontal slice of $\mathcal{A}$ in the frequency domain represents a classical Hermitian positive semidefinite form, hence $\mathcal{A}$ represents a $t$-Hermitian positive semidefinite form. Furthermore, by equation \eqref{contraction hypermatrix in frequency domain}, we have that  
    \[\Phi_c(\mathcal{A}) = \frac{1}{p}\overline{\widehat{c}[m]}\mathcal{B},\]
    thus the $c$-contraction hypermatrix $\Phi_c(\mathcal{A})$ represents the classical Hermitian positive semidefinite form
    \[h_{\Phi_c(\mathcal{A})}(x) = \frac{1}{p}\overline{\widehat{c}[m]}h_{\mathcal{B}}(x).\]
    Moreover, since $\mathcal{B}\neq 0$ and is Hermitian positive semidefinite, this implies that there exists $z_*\in \mathbb{C}^n$ such that 
    \[h_{\mathcal{B}}(z_*)>0.\]
    Therefore, Hermitian positive semidefiniteness of $\frac{1}{p}\overline{\widehat{c}[m]}h_{\mathcal{B}}(x)$ implies that 
    \[\frac{1}{p}\overline{\widehat{c}_m}h_{\mathcal{B}}(z_*)\geq 0,\]
    hence $\overline{\widehat{c}[m]}\geq 0$; that is, $\widehat{c}[m]\geq 0$. Since $m$ was arbitrary, repeatedly applying this construction yields $\widehat{c}[\ell]\geq 0$ for all $l\in [p]$. Thus, $(ii)\Rightarrow (i)$. 

    Next, suppose $c\in \mathbb{T}_p^+\setminus \{0\}$ and let $\mathcal{A}\in \mathbb{C}^{n\times\dots\times n\times p}$ be any order $2k+1$ $t$-Hermitian positive definite hypermatrix. By $t$-Hermitian positive definiteness, for every nonzero $z\in \mathbb{C}^n$ we have that 
    \[h_{\widehat{\mathcal{A}}_{(\ell)}}(z) > 0\]
    for each $\ell\in [p]$. Consequently, 
    \[h_{\Phi_c(\mathcal{A})}(z) = \frac{1}{p}\sum\limits_{\ell=1}^p\widehat{c}[\ell]h_{\widehat{\mathcal{A}}_{(\ell)}}(z)>0\]
    for all nonzero $z\in \mathbb{C}^n$, proving that classical Hermitian form induced by $\Phi_c(\mathcal{A})$ is Hermitian positive definite. Thus, $(a)\Rightarrow (b)$. Conversely, suppose $\Phi_c$ strongly preserves positivity; that is, for every order $2k+1$ $t$-Hermitian positive definite hypermatrix $\mathcal{A}\in \mathbb{C}^{n\times\dots\times n\times p}$, $h_{\Phi_c(\mathcal{A})}(x)$ is Hermitian positive semidefinite. For any order $2k+1$ $t$-Hermitian positive definite hypermatrix $\mathcal{A}\in \mathbb{C}^{n\times\dots\times n\times p}$, let $\epsilon > 0$ and $\mathcal{B}\in \mathbb{C}^{n\times\dots\times n\times p}$ be an order $2k+1$ $t$-Hermitian positive definite hypermatrix. Then $\mathcal{A}+\epsilon\mathcal{B}$ is $t$-Hermitian positive definite, and by assumption 
    \[\Phi_c(\mathcal{A}+\epsilon\mathcal{B}) = \Phi_c(\mathcal{A})+\epsilon\Phi_c(\mathcal{B})\]
    represents a strictly positive Hermitian form. That is, for any nonzero $z\in \mathbb{C}^n$, 
    \[h_{\Phi_c(\mathcal{A})}(z)+\epsilon h_{\Phi_c(\mathcal{B})}(z) > 0,\]
    and so taking $\epsilon\rightarrow 0^+$, it follows that 
    \[h_{\Phi_c(\mathcal{A})}(z)\geq 0.\]
    Hence, $\Phi_c(\mathcal{A})$ is Hermitian positive semidefinite, showing that $\Phi_c$ also weakly preserves positivity. From the previous paragraph, this implies that $c\in \mathbb{T}_p^+$. Indeed, it must be the case that $c\in \mathbb{T}_p^+\setminus \{0\}$, for otherwise if $c=0$, then $\Phi_c(\mathcal{A})=0$ for all $\mathcal{A}\in \mathbb{C}^{n\times\dots\times n\times p}$, which does not represent a Hermitian positive definite form, violating the assumption that $\Phi_c$ strongly preserves positivity. Thus, $(b)\Rightarrow (a)$. 
\end{proof}

\end{document}
